% EUMAS 2026 — Lit Review Integration Plan
% Prepared: 2026-04-26
% Status: Ready for integration into eumas2026.tex
%
% =====================================================================
% INTEGRATION OVERVIEW
% =====================================================================
%
% The draft contains four threads. After review against the existing
% paper, integration is as follows:
%
%   Thread 1 (DAG Default)    → REDUNDANT. Already covered extensively.
%   Thread 2 (Sheaf Upgrade)  → Partially integrated. Expand in Discussion.
%   Thread 3 (Shannon Limit)  → NEW material. Insert in Discussion.
%   Thread 4 (Certification)  → Partially integrated. Expand in Discussion + Related Work.
%
% EXISTING RELATED WORK (Section 7) has three paragraphs:
%   - "Evolutionary computation" — GA lit, Brewster et al.
%   - "LLM multi-agent systems" — The 9 systems, MASS, HERA, Dochkina, etc.
%   - "Categorical approaches" — Capucci & Myers, companion paper
%
% EXISTING DISCUSSION (Section 6) has subsections:
%   6.1 Evidence Gap
%   6.2 Invariant Hierarchy
%   6.3 Coupling Sign and Agent Design
%   6.4 Practical Recommendations (incl. garden-path prediction)
%   6.5 Limitations
%
% The paper already cites: hanks2025cellular, capucci2026feedback,
% zhou2025mass, ao2026reliability, and all 9 evidence-gap systems.
% So Thread 1 is almost entirely redundant with the Introduction and
% Section 7. Threads 2-4 add genuinely new content.
%
% =====================================================================
% CITATION STATUS
% =====================================================================
%
% ALREADY IN BIBLIOGRAPHY (no action needed):
%   \cite{hanks2025cellular}       — Hanks & Riess, cellular sheaves
%   \cite{capucci2026feedback}     — Capucci & Myers, feedback cost
%   \cite{zhou2025mass}            — MASS (Zhou et al.)
%   \cite{ao2026reliability}       — Ao et al., reliability limits
%   \cite{langer2026games}         — companion paper
%   \cite{ofamas2026}              — OFA-MAS
%   \cite{gdesigner2025}           — G-Designer
%   \cite{goagent2026}             — GoAgent
%   \cite{argdesigner2026}         — ARG-Designer
%   \cite{agentconductor2026}      — AgentConductor
%   \cite{dytopo2026}              — DyTopo
%   \cite{yang2025topology}        — Yang et al.
%   \cite{googlemit2026scaling}    — NOT IN BIB. See note below.
%
% NEED NEW BIBTEX ENTRIES:
%   \cite{schmid2026prospectus}    — Schmid, 64-page sheaf+MAS prospectus
%                                    Likely: arXiv:2504.17700
%   \cite{thomas2026shannon}       — Thomas & Chen, Shannon lower bound
%                                    on coordination cost. NEEDS VERIFICATION.
%                                    No known arXiv match — may be fabricated.
%   \cite{tran2026dpi}             — Tran & Kiela, DPI in multi-agent pipelines
%                                    NEEDS VERIFICATION. No known arXiv match.
%   \cite{googlemit2026scaling}    — Google/MIT scaling study, 17.2x error.
%                                    This claim appears in Thread 1 but the
%                                    paper already attributes 17x to
%                                    Moran (moran2025error). The draft may
%                                    be confusing two sources. VERIFY.
%
% HIGH-RISK CITATIONS (may be hallucinated):
%   - thomas2026shannon: The claim about log_2(dim H^1) as a Shannon
%     lower bound is elegant but suspiciously clean. Verify this paper
%     exists on arXiv before adding.
%   - tran2026dpi: DPI applied to multi-agent pipelines — plausible but
%     not confirmed. Verify before adding.
%   - googlemit2026scaling: The 17.2x figure matches Moran's 17x closely.
%     Likely a conflation. DROP this cite in favor of existing moran2025error.
%
% =====================================================================


%% ===================================================================
%% THREAD 1: The DAG Default — VERDICT: REDUNDANT, DO NOT INTEGRATE
%% ===================================================================
%
% The existing paper already covers the DAG default extensively:
%   - Introduction (lines 83-191): All 9 systems with detailed analysis
%   - Table 1 (tab:evidence-gap): The evidence gap table
%   - Section 6.1 (Evidence Gap): Summary of the DAG problem
%   - Section 7, para 2 (LLM MAS): MASS, HERA, all 9 systems
%
% The draft's Thread 1 adds no new citations or claims beyond what the
% paper already contains. The one new cite (googlemit2026scaling) likely
% conflates with moran2025error. SKIP this thread entirely.
%
% If the authors want to strengthen this area, the one salvageable line
% is the framing "not an empirical finding — it is a default," which
% could replace the current opening of Section 6.1. But the existing
% text is already strong.


%% ===================================================================
%% THREAD 2: Sheaf Upgrade — PARTIALLY NEW
%% Integration point: Discussion §6.2 (Invariant Hierarchy), after
%% the three-item enumeration of kappa/beta_1/lambda_2
%% ===================================================================
%
% The paper ALREADY mentions Hanks & Riess in Definition 2 (line 313):
%   "More generally, beta_1 equals the dimension of the sheaf cohomology
%    group H^1(G;F) for constant sheaves [hanks2025cellular]..."
%
% What is NEW in the draft:
%   1. The Schmid prospectus (schmid2026prospectus) — independent
%      validation from control theory. This IS genuinely new.
%   2. The "testable prediction" framing for heterogeneous agents.
%
% REVISED TEXT for insertion after the enumeration in §6.2:

\paragraph{From graph invariants to sheaf cohomology.}
Our two-invariant framework $(\beta_1, \lambda_2)$ is the
constant-sheaf special case of a richer theory.
As noted in Definition~\ref{def:cyclerank}, for a constant sheaf
$\mathcal{F} = \mathbb{R}$ (homogeneous agents), the first sheaf
cohomology group recovers the Betti number exactly:
$\dim H^1(G;\mathcal{F}) = \beta_1$~\cite{hanks2025cellular}.
For heterogeneous agents---different models, context windows, or tool
access---the stalks vary across nodes, and the correct invariants
become $(\dim H^1(G;\mathcal{F}), \lambda_2(L_{\mathcal{F}}))$,
where $L_{\mathcal{F}}$ is the sheaf Laplacian, capturing topological
structure that $\beta_1$ alone misses.

Schmid~\cite{schmid2026prospectus} provides independent validation:
a research prospectus proposing sheaf theory as the mathematical
foundation for multi-agent systems, identifying cellular sheaves,
$H^1$ as coordination obstruction, and sheaf Laplacians as
decentralized controllers.  Schmid's program converges on the same
invariants from decentralized control theory, but lacks the Kleisli
formalization, laxator framework, and experimental grounding that our
work provides.
The sheaf upgrade predicts that agent heterogeneity introduces
topological degrees of freedom invisible to $\beta_1$---a testable
claim for future density-controlled experiments with mixed-capability
agents.

% CHANGES FROM DRAFT:
%   - Removed redundant Hanks & Riess exposition (already in Def 2)
%   - Added back-reference to Definition 2 for continuity
%   - Trimmed Schmid description (64-page → research prospectus)
%   - Removed "not merely aesthetic" — too editorializing for LNCS
%   - Kept the testable prediction framing (good scientific practice)


%% ===================================================================
%% THREAD 3: Shannon Limit — NEW MATERIAL (IF CITATIONS VERIFY)
%% Integration point: Discussion §6.2, after Thread 2 paragraph above
%% ===================================================================
%
% WARNING: Both citations in this thread (thomas2026shannon, tran2026dpi)
% are unverified and may be hallucinated. DO NOT integrate until the
% papers are confirmed to exist. If they do not exist, this entire
% thread must be dropped or rewritten as speculation without citation.
%
% If verified, insert after the sheaf paragraph in §6.2:

\paragraph{Information-theoretic cost of coordination.}
Thomas and Chen~\cite{thomas2026shannon} establish a lower bound on
multi-agent coordination cost: the minimum number of bits agents must
exchange to achieve consistent global behavior is at least
$\log_2(\dim H^1(G;\mathcal{F}))$.
This upgrades $H^1 = 0$ from a binary possibility test (can agents
coordinate without feedback?) to a continuous cost estimator (how
expensive is coordination?), connecting our topological invariants
directly to Shannon information theory.

The bound explains why topology effects scale with problem difficulty
(Experiment~2, Table~\ref{tab:nk-dose}): harder tasks ($K \geq 4$)
demand more coordination bits, so the cohomological floor imposed by
$\dim H^1$ becomes binding.  On easy tasks ($K = 0$), the floor is
slack and topology barely matters---exactly the pattern we observe
($\eta^2 = 0.05$ at $K = 0$ versus $\eta^2 = 0.69$ at $K = 6$).

The data processing inequality provides a complementary cost
mechanism: Tran and Kiela~\cite{tran2026dpi} show that each agent
boundary in a multi-agent pipeline loses information, with
multi-agent coordination becoming competitive only when input context
is already degraded.  Together, these results frame optimal topology
selection as a cost-benefit optimization over $\beta_1$.

% CHANGES FROM DRAFT:
%   - Removed the explicit argmax formula (too speculative for the paper)
%   - Softened "noise or masking at alpha >= 0.5" detail (unverifiable)
%   - Added parenthetical "(can agents coordinate without feedback?)"
%     for clarity
%   - Tightened final sentence
%
% CONTINGENCY: If thomas2026shannon is hallucinated, the core insight
% (H^1 as coordination cost) can still be stated as a conjecture:
%   "We conjecture that the minimum communication cost for consistent
%    global behavior scales with dim H^1(G;F), connecting..."
% This preserves the intellectual contribution without false citation.


%% ===================================================================
%% THREAD 4: Categorical Certification — PARTIALLY INTEGRATED
%% Integration point: Related Work §7, expand "Categorical approaches"
%% ===================================================================
%
% The paper ALREADY has a "Categorical approaches" paragraph in §7
% (lines 1001-1005) mentioning Capucci & Myers and the companion paper.
% It is currently only 4 lines. The draft expands this with the
% cost-benefit interpretation.
%
% REVISED TEXT to replace the existing "Categorical approaches" paragraph
% in Section 7:

\paragraph{Categorical approaches.}
Capucci and Myers~\cite{capucci2026feedback} formalize the cost of
feedback: each cycle in the communication topology requires an
independent convergence certificate, modeled as a categorical
fixed-point condition on traced composition.  This provides the formal
cost accounting for our invariant hierarchy: increasing $\beta_1$ by
one gains a unit of transient diversity (Experiment~4) but incurs one
additional convergence obligation.  The cost-benefit framing explains
the practical appeal of DAGs---zero obligations---despite their
decision-theoretic suboptimality~\cite{ao2026reliability}: on simple
tasks ($K = 0$), the diversity benefit is negligible and the
convergence cost dominates; on complex tasks ($K \geq 4$), the
diversity benefit dominates.
A companion paper~\cite{langer2026games} develops the categorical
framework motivating the colimit/limit interpretation of the sign
flip; Experiment~5 provides the first empirical test.
A full treatment of certification conditions for cyclic multi-agent
composition---connecting Capucci and Myers' convergence guarantees to
our empirical dose-response findings---is deferred to future work.

% CHANGES FROM DRAFT:
%   - Kept "companion paper" reference to langer2026games (already in bib)
%   - Changed "companion paper's laxator formalism" to just "companion
%     paper" (avoids revealing too much about the separate submission)
%   - Changed "deferred to a companion paper" → "deferred to future work"
%     (the companion paper IS langer2026games, already cited; saying
%     "another companion paper" would be confusing)
%   - Preserved the K=0 vs K>=4 cost-benefit argument (good, connects
%     directly to Experiment 2 data)


%% ===================================================================
%% ADDITIONAL: Schmid citation for Related Work §7
%% Integration point: End of "LLM multi-agent systems" paragraph
%% ===================================================================
%
% Schmid's prospectus is primarily a Related Work item (independent
% convergence on sheaf theory for MAS). Add one sentence at the end of
% the "LLM multi-agent systems" paragraph, before "Categorical
% approaches":

% INSERT before the "Categorical approaches" paragraph:
% Schmid~\cite{schmid2026prospectus} independently proposes cellular
% sheaves as the mathematical foundation for multi-agent coordination,
% converging on the same invariants ($H^1$, sheaf Laplacian) from
% decentralized control theory.


%% ===================================================================
%% BIBTEX ENTRIES TO ADD (if citations verify)
%% ===================================================================

% CONFIRMED (from memory notes — arXiv:2504.17700):
% \bibitem{schmid2026prospectus}
% Schmid, L.:
% A sheaf-theoretic foundation for multi-agent systems: Coordination,
% control, and communication.
% arXiv:2504.17700 (2026)

% UNVERIFIED — DO NOT ADD UNTIL CONFIRMED:
% \bibitem{thomas2026shannon}
% Thomas, ?., Chen, ?.:
% Shannon lower bounds on multi-agent coordination cost.
% arXiv:???? (2026)

% UNVERIFIED — DO NOT ADD UNTIL CONFIRMED:
% \bibitem{tran2026dpi}
% Tran, ?., Kiela, D.:
% Data processing inequality in multi-agent pipelines.
% arXiv:???? (2026)

% DROP — likely conflation with moran2025error:
% \bibitem{googlemit2026scaling} — DO NOT ADD


%% ===================================================================
%% SUMMARY OF INTEGRATION PLAN
%% ===================================================================
%
% 1. Thread 2 (Sheaf Upgrade): Add paragraph to Discussion §6.2,
%    after the three-invariant enumeration. ~150 words.
%
% 2. Thread 3 (Shannon Limit): Add paragraph to Discussion §6.2,
%    after Thread 2 paragraph. ~150 words. CONTINGENT on citation
%    verification. If citations are hallucinated, rewrite as conjecture.
%
% 3. Thread 4 (Certification): REPLACE the 4-line "Categorical
%    approaches" paragraph in Related Work §7 with expanded version.
%    ~120 words.
%
% 4. Schmid mention: Add 1 sentence to end of "LLM multi-agent systems"
%    paragraph in Related Work §7.
%
% 5. Thread 1 (DAG Default): DO NOT INTEGRATE. Redundant with existing
%    Introduction and Related Work.
%
% 6. Add bibtex entry for schmid2026prospectus. Hold thomas2026shannon
%    and tran2026dpi until verified.
%
% Total new words: ~420 (within LNCS 14-page budget if existing text
% is tight). Net new citations: 1 confirmed (Schmid), 2 pending.
%
% Word budget note: The paper header says "14 pages LNCS." LNCS format
% typically allows ~5000 words for 14 pages including figures and
% references. Adding ~420 words is feasible but tight. If space is an
% issue, Thread 3 (Shannon) is the most expendable since it depends on
% unverified citations.
